\begin{textsample}{POS Dim 2 – 2018\_c4 – Score 16.00 – 2018\_c4/enade\_2018\_psico\_860b383a75.txt}  \label{ex:f2_pos_056}
CAPÍTULO 4 PERCEPÇÃO DA PROVA Nas análises feitas neste capítulo, constam as percepções dos concluintes da Área de Psicologia sobre a prova aplicada no Enade/2018.
Essas percepções foram mensuradas por meio de nove questões que avaliaram, desde o grau de dificuldade da prova até o tempo gasto para concluí -la.
As percepções sobre a prova foram relacionadas ao desempenho dos estudantes e à Grande Região de funcionamento do curso.
O \textbf{questionário} de percepção da prova encontra -se ao final do Anexo VII, que traz a \textbf{reprodução} do exame.
O desempenho dos estudantes foi classificado em quatro quartos.
Para tanto, esse desempenho foi \textbf{ordenado} de forma ascendente.
O percentil 25, P25, também conhecido como primeiro quartil, é a nota de desempenho que deixa um quarto ( 25 \% ) dos valores observados abaixo e três quartos acima.
A Figura 1 apresenta uma ilustração desse conceito.
O quarto inferior de desempenho é composto pelas notas abaixo do primeiro quartil.
Já o percentil 75, P75, também conhecido como terceiro quartil, é o valor para o qual há três quartos ( 75 \% ) dos dados abaixo e um quarto acima dele.
O quarto superior de desempenho é composto pelas notas iguais ou acima do terceiro quartil.
O percentil 50, P50, também conhecido como mediana, é o valor que divide as notas em dois conjuntos de igual tamanho.
O segundo quarto inclui valores entre o primeiro quartil ( P25 ) e a mediana.
O terceiro quarto contém os valores entre a mediana ( P50 ) e o terceiro quartil ( P75 ).
Vale ressaltar que percentis, quartis e medianas são pontos que não, obrigatoriamente, pertencem ao conjunto original de dados, ao passo que os quartos são \textbf{subconjuntos} dos dados originais.
A seguir, serão apresentados gráficos com resultados selecionados, relativos às nove questões avaliadas por grupos de estudantes.
Nas barras dos gráficos, apresenta -se o percentual de alunos que assinalaram uma das opções ou a soma das porcentagens daqueles que assinalaram duas ( ou três ) delas.
Por exemplo, para as questões 1 e 2, nos gráficos, é apresentada a porcentagem total de participantes que assinalaram as opções ( D ) difícil e ( E ) Muito difícil.
Em cada barra, foram assinalados, também os extremos do intervalo de confiança de 95 \% como linhas verticais unidas por uma linha horizontal na forma da letra H maiúscula.
O \textbf{estimador} de um \textbf{parâmetro} com um certo nível de confiança ( e.g. 95,0 \% ) deve conter o \textbf{parâmetro} no intervalo de confiança em 95 \% das vezes.
Na comparação entre os \textbf{estimadores} dos \textbf{parâmetros} de duas classes de uma dada categoria ( e.g.
Norte e Nordeste nas Grandes Regiões ou de primeiro e último quarto dentro de desempenho ) associados aos seus respectivos intervalos de confiança, diz -se que não há diferença estatisticamente significativa entre os \textbf{parâmetros} das duas categorias quando há \textbf{interseção} entre os intervalos de confiança e que há diferença, se os intervalos de confiança são \textbf{disjuntos} ( para maiores informações \textbf{vide} Glossário ).
Nas Tabelas no Anexo II, são apresentados os valores absolutos e a distribuição percentual18 das alternativas válidas das nove questões, segundo o mesmo \textbf{recorte} de desempenho dos alunos e Grande Região de funcionamento do curso. 4.1 GRAU DE DIFICULDADE DA PROVA 4.1.1 Componente de Formação Geral Ao avaliarem Qual o grau de dificuldade desta prova na parte de Formação Geral? ( Questão 1 ), 28,2 \% do grupo de inscritos e presentes optaram pelas alternativas Difícil ou Muito difícil.
Entretanto, para mais da metade dos estudantes ( 61,4 \% ), o Componente de Formação Geral da prova foi considerado com grau de dificuldade Médio ( Gráfico 4.1, Gráfico 4.2 e, no Anexo II, a Tabela II.1 ).
O percentual de estudantes que consideraram a prova Difícil ou Muito difícil foi maior na região Sul, onde a proporção foi de 32,2 \%, enquanto a de menor incidência foi a Nordeste, com 24,0 \%.
No Gráfico 4.1, é possível observar que as diferenças entre as regiões Nordeste e Sudeste com as demais regiões são estatisticamente significativas.
Nas Grandes Regiões, a proporção de presentes à prova que consideraram o Componente de Formação Geral como sendo de grau de dificuldade Médio esteve entre 58,1 \%, na região Sul, e 65,4 \%, na região Nordeste. 18 Cumpre lembrar uma das convenções para tabelas \textbf{numéricas} ( pág. iii ) sobre a possibilidade de a soma das partes não resultar em 100 \% por questões de \textbf{arredondamento}. sendo de grau de dificuldade Médio esteve entre 58,1 \%, na região Sul, e 65,4 \%, na região Nordeste.
O percentual de alunos que consideraram a prova Difícil ou Muito difícil foi decrescente em função dos quartos de desempenho : 35,1 \%, no primeiro quarto e 18,6 \%, no quarto superior, grupo de melhor desempenho na prova.
Nos quartos de desempenho intermediários, a proporção de alunos que consideraram a prova Difícil ou Muito difícil foi de 32,1 \%, no segundo quarto e de 27,6 \%, no terceiro quarto.
As diferenças entre os quartos de desempenho são estatisticamente significativas.
Para todos os quartos de desempenho, a alternativa modal para essa pergunta foi a Médio, com 58,2 \% e 63,9 \% dos respondentes nos quartos extremos, primeiro e quarto, respectivamente. 4.1.2 Componente de Conhecimento Específico Ao responderem à Questão 2 – Qual o grau de dificuldade desta prova na parte de Componente Específico? – 30,2 \% do grupo de estudantes classificaram -na como Difícil ou Muito difícil.
Além disso, o Componente de Conhecimento Específico da prova foi considerado com grau de dificuldade Médio por 62,8 \% dos alunos ( Gráfico 4.3, Gráfico 4.4, e, no Anexo II, a Tabela II.2 ).
A análise das respostas dos estudantes quanto ao grau de dificuldade do Componente de Conhecimento Específico da prova, agregado por Grande Região, mostra que a diferença entre a menor e as duas maiores proporções de alunos que a avaliaram como Difícil ou Muito difícil é estatisticamente significativa : a menor na região Nordeste ( 24,7 \% ) e as duas maiores proporções nas regiões Sul ( 34,5 \% ) e Centro-Oeste ( 33,8 \% ).
O percentual de alunos que classificaram o grau de dificuldade como Médio, no Componente de Conhecimento Específico, variou de 59,1 \% a 67,0 \%, para as regiões Sul e Nordeste, respectivamente.
Considerando -se a avaliação da dificuldade das questões do Componente de Conhecimento Específico da prova, de acordo com o desempenho dos estudantes, observa -se que há diferença estatisticamente significativa dos resultados entre o quarto superior de desempenho e os demais quartos.
A proporção dos que classificaram a parte específica como Difícil ou Muito difícil variou de 27,3 \% ( último quarto ) a 31,5 \% ( segundo quarto ).
As demais proporções dessa reposta foram de 32,6 \% e 29,7 \%, no primeiro e terceiro quarto, respectivamente.
Já a proporção dos que responderam que o grau de dificuldade das questões do Componente de Conhecimento Específico da prova foi Médio variou de 60,4 \%, no primeiro quarto a 64,4 \%, no quarto superior. 4.2 EXTENSÃO DA PROVA EM RELAÇÃO AO TEMPO TOTAL \textbf{Indagados} quanto à extensão da prova, em relação ao tempo total oferecido para a sua resolução ( Questão 3 ), os estudantes apontaram, com maior incidência, a alternativa que considerava a extensão adequada, para todas as agregações consideradas ( Gráfico 4.5, Gráfico 4.6, e, no Anexo II, a Tabela II.3 ).
O percentual de alunos que responderam ser a extensão da prova adequada foi de 45,2 \%.
Já 46,6 \% dos inscritos presentes consideraram que a prova foi longa ou muito longa, e 8,1 \% a avaliaram como curta ou muito curta.
Dentre as Grandes Regiões, a proporção daqueles que avaliaram a prova como longa ou muito longa, em relação ao tempo total destinado à sua resolução, variou de 41,4 \%, na região Norte até 48,1 \%, na região Sul.
Há diferença estatisticamente significativa entre a região Norte e as demais regiões.
Considerando -se o desempenho dos alunos, nota -se, ainda, que 42,8 \% consideraram a extensão da prova adequada no quarto de desempenho inferior e 48,5 \% no de melhor desempenho ( quarto superior ).
Nos quartos intermediários, essa proporção foi de 43,4 \%, no segundo quarto e de 46,1 \%, no terceiro.
No Gráfico 4.6, pode -se constatar que há diferença estatisticamente significativa dentre as proporções de estudantes que consideraram a prova longa ou muito longa em função dos quartos de desempenho entre o último quarto ( melhor desempenho ) com o primeiro e o segundo quarto, sendo a maior proporção dessa resposta no primeiro quarto de desempenho ( 47,8 \% ), e a menor, no quarto inferior ( 45,1 \% ). 4.3 COMPREENSÃO DOS ENUNCIADOS DAS QUESTÕES 4.3.1 Componente de Formação Geral Em relação aos enunciados das questões do Componente de Formação Geral ( Questão 4 ), as opiniões foram positivas, já que 80,0 \% dos alunos avaliados consideraram todos ou a maioria dos enunciados das questões claros e objetivos ( Gráfico 4.7, Gráfico 4.8, e, no Anexo II, a Tabela II.4 ).
Na análise regional, a \textbf{percentagem} de estudantes que avaliaram que todos ou a maioria dos enunciados das questões do Componente de Formação Geral estavam claros e objetivos variou de 76,8 \%, na região Nordeste a 83,1 \%, na região Sul.
A análise das percepções dos estudantes sobre a clareza e objetividade dos enunciados permite afirmar que todos, ou a maioria dos enunciados de questões relativas ao Componente de Formação Geral, foram considerados claros e objetivos para a maior parte dos respondentes ( maior ou igual a 76,8 \% em todas as regiões e maior ou igual a 69,7 \% para todos os quartos de desempenho ).
Segundo o desempenho, observa -se que a proporção dos que emitiram essa opinião \textbf{cresce} conforme o desempenho aumenta, com diferenças estatisticamente significativas entre os quartos de desempenho.
No quarto superior, a clareza e objetividade de todos ou da maioria dos enunciados das questões foram percebidas por 88,8 \%. 4.3.2 Componente de Conhecimento Específico Em relação aos enunciados das questões do Componente de Conhecimento Específico da prova, para 81,6 \% dos estudantes avaliados da Área de Psicologia, a clareza e a objetividade ( Questão 5 ) estavam presentes em todas ou na maioria das questões ( Gráfico 4.9, Gráfico 4.10, e no Anexo II, a Tabela II.5 ).
A maioria dos estudantes de todas as Grandes Regiões brasileiras considerou claros e objetivos todos ou a maioria dos enunciados das questões do Componente de Conhecimento Específico da prova, percentual sempre maior ou igual a 80,2 \%.
As diferenças entre a região Sul com as regiões Nordeste e Sudeste, além da entre a região Centro-Oeste e a região Nordeste, são estatisticamente significativas.
A proporção de estudantes que consideraram os enunciados das questões claros e objetivos apresenta uma tendência crescente com o aumento de desempenho : mais elevada no quarto superior ( 87,6 \% ), se comparada ao quarto inferior de desempenho ( 73,3 \% ).
As diferenças entre o todos os quartos de desempenho são estatisticamente significativas. 4.4 SUFICIÊNCIA DAS INFORMAÇÕES/INSTRUÇÕES FORNECIDAS Ao avaliarem as informações/instruções fornecidas para a resolução das questões ( Questão 6 ), 88,0 \% dos respondentes da Área de Psicologia de todo o Brasil afirmaram que estas eram até excessivas ou suficientes em todas ou na maioria das questões ( Gráfico 4.11, Gráfico 4.12, e, no Anexo II, a Tabela II.6 ).
Quanto à distribuição de respondentes pelas Grandes Regiões, observa -se que a proporção de estudantes que consideraram as informações/instruções fornecidas até excessivas ou suficientes em todas ou na maioria das questões foi sempre superior ou igual a 86,1 \%, chegando a 89,5 \%, na região Sul.
Não há diferença estatisticamente significativa entre as regiões Norte e Nordeste, nem entre as regiões Sudeste e Sul e nem entre a região Centro-Oeste com as outras regiões.
Levando -se em conta o desempenho dos participantes, observa -se uma tendência crescente das proporções de participantes que avaliaram as informações/instruções como até excessivas ou suficientes em todas ou na maioria das questões, havendo diferenças estatisticamente significativas entre todos os quartos de desempenho, como mostrado no Gráfico 4.12.
O percentual foi mais elevado no quarto superior ( 93,0 \% ), acima da média nacional ( 88,0 \% ).
Já no quarto inferior, a suficiência das informações/instruções declarada como até excessiva, em todas ou na maioria das questões foi percebida por 82,0 \% dos respondentes. 4.5 DIFICULDADE ENCONTRADA AO RESPONDER À PROVA \textbf{Indagados} sobre as dificuldades com as quais se depararam ao responder à prova ( Questão 7 ), 15,9 \% dos estudantes apontaram o Desconhecimento do conteúdo.
Para 41,7 \%, a Forma diferente de abordagem do conteúdo foi indicada como dificuldade.
Já a Falta de motivação para fazer a prova foi a dificuldade apontada por 15,5 \% dos respondentes.
Considerando -se todo o Brasil, 20,5 \% dos respondentes afirmaram que não tiveram qualquer tipo de dificuldade para responder à prova ( Tabela II.7 no Anexo II ).
Nos Gráficos 4.13 e 4.14, são apresentados os percentuais de estudantes que apontaram o Desconhecimento do conteúdo como dificuldade ao responder à prova.
Na análise por Grandes Regiões, o percentual de inscritos e presentes que apontaram o Desconhecimento do conteúdo como dificuldade ao responder à prova não superou 17,4 \%.
Os percentuais variaram de 13,9 \%, na região Nordeste, a 17,4 \%, na Sul.
A Forma diferente de abordagem do conteúdo foi a escolha modal dos estudantes, com percentuais que variaram de 39,5 \% ( região Nordeste ) a 43,1 \% ( região Centro-Oeste ).
O percentual de alunos que citaram a Falta de motivação como dificuldade variou de 12,2 \% ( região Norte ) a 17,8 \% ( região Nordeste ).
Os que declararam não ter qualquer tipo de dificuldade para responder à prova variaram de 19,0 \%, na região Sul, a 21,1 \%, na região Sudeste.
Em relação aos quartos de desempenho, o Desconhecimento do conteúdo foi a opção escolhida por 15,8 \% dos estudantes do quarto superior e por 16,0 \% do quarto inferior, sendo esse o mesmo valor para os quartos intermediários.
Foram observadas a não existência de diferenças estatisticamente significativas entre as proporções dos quartos de desempenho.
A alternativa modal para os alunos, quando agregados pelos quartos de desempenho, foi que a dificuldade encontrada foi causada pela Forma diferente de abordagem do conteúdo : 43,7 \%, no segundo quarto e 39,6 \% no quarto superior escolheram essa alternativa. 4.6 CONTEÚDOS DAS QUESTÕES OBJETIVAS DA PROVA Ao analisar os conteúdos das questões objetivas da prova ( Questão 8 ), um percentual muito pequeno dos estudantes avaliados, apenas 3,2 \%, afirmou que não estudou ainda a maioria desses conteúdos ( Gráficos 4.15, Gráfico 4.16, e a Tabelas II.8 no Anexo II ).
A maioria ( 78,4 \% ) afirmou ter estudado e aprendido muitos ou todos os conteúdos.
Na análise por Grande Região, a proporção de respondentes que escolheram a opção não estudou ainda a maioria desses conteúdos foi pequena.
Observa -se que, nas regiões Norte ( 4,9 \% ) e Nordeste ( 3,7 \% ), as proporções foram maiores do que a média nacional ( 3,2 \% ).
Foram observadas diferenças estatisticamente significativas entre a região Norte e as regiões Sudeste, Sul e Centro-Oeste.
Também foi observada diferença estatisticamente significativa entre as regiões Nordeste e Sul.
Em todas as regiões, a maioria dos presentes afirmou ter estudado e aprendido muitos ou todos os conteúdos, com proporções variando entre 70,6 \%, na região Norte, a 79,3 \%, na região Sudeste.
Considerando -se, separadamente, as opiniões de estudantes dos quartos de desempenho, observa -se que, no quarto inferior, 6,4 \% ofereceram como resposta que não estudou ainda a maioria desses conteúdos, sendo 1,1 \% os do quarto superior com a mesma resposta.
As diferenças, entre os alunos que optaram por esse motivo de dificuldade em todos os quartos de desempenho são estatisticamente significativas.
Tendo -se em conta o quarto superior, 88,4 \% dos alunos afirmaram ter estudado e aprendido muitos ou todos os conteúdos.
No outro extremo, no primeiro quarto, 65,9 \% optaram pela mesma categoria.
As proporções são crescentes com o desempenho. 4.7 TEMPO GASTO PARA CONCLUIR A PROVA Ao responderem sobre o tempo de conclusão da prova ( Questão 9 ), um pouco mais de três quartos dos estudantes ( 80,1 \% ) afirmaram ter gasto Entre duas e quatro horas ( Gráfico 4.17, Gráfico 4.18 e, no Anexo II, a Tabela II.9 ).
Considerando -se as cinco Grandes Regiões brasileiras, os percentuais dos que utilizaram Entre duas e quatro horas para finalizar a prova nas regiões Norte ( 77,2 \% ), Nordeste ( 77,6 \% ) e Centro-Oeste ( 78,4 \% ) foram inferiores ao percentual nacional.
Nas regiões Sudeste e Sul o percentual de alunos que \textbf{dispensaram} Entre duas e quatro horas para concluir a prova ficou igual ou acima de 81,1 \%, como mostrado no Gráfico 4.17.
Há diferenças estatisticamente significativas entre as regiões que ficaram acima do percentual nacional com as que ficaram abaixo.
Uma vez analisadas as alternativas escolhidas pelos estudantes que se situam nos diferentes quartos de desempenho.
Observa -se uma tendência crescente : uma proporção maior de participantes no quarto superior declarou ter gasto Entre duas e quatro horas para concluir a prova quando comparada com a proporção de participantes nos quartos inferiores.
Há diferenças estatisticamente significativas entre os quartos de desempenho, à exceção do terceiro com o último quarto.

% matched lemmas: arredondamento, crescer, disjunto, dispensar, estimador, indagar, interseção, numérico, ordenar, parâmetro, percentagem, questionário, recorte, reprodução, subconjunto, vidar
\end{textsample}
